% --- Содержимое файла materials/2_1.tex ---
\subsection{Теорема о вложенных шарах и теорема Бэра}
\begin{definition}[Полное метрическое пространство]
	МП \textit{полное}, если любая фундаментальная последовательность сходится.
\end{definition}

\begin{theorem}[О вложенных шарах]\label{kantor-compact}
	\textbf{Условия:}
	\begin{itemize}
		\item $X$ --- полное метрическое пространство;
		\item $\{\overline{B}_{r_n}(x_n)\}$ --- последовательность вложенных замкнутых шаров ($\overline{B}_{r_{n+1}} \subset \overline{B}_{r_n}$);
		\item $r_n \to 0$ при $n \to \infty$.
	\end{itemize}
	\textbf{Следовательно:}
	\begin{itemize}
		\item $\bigcap\limits_{n \in \mathbb{N}} \overline{B}_{r_n}(x_n) \neq \varnothing$;
		\item Пересечение состоит ровно из одной точки: $\bigcap\limits_{n \in \mathbb{N}} \overline B_{r_n}(x_n) = \{x^*\}$.
	\end{itemize}
\end{theorem}

\begin{proof}
	\begin{enumerate}
		\item \textbf{Фундаментальность:} В силу вложенности шаров $\overline{B}_{r_{n+p}} \subset \overline{B}_{r_n}$ расстояние между центрами $\rho(x_n, x_{n+p}) \leqslant r_n$. Так как по условию $r_n \to 0$, последовательность $\{x_n\}$ является фундаментальной.
		\item \textbf{Сходимость:} Поскольку пространство $X$ полно по определению, существует точка $x^* \in X$, такая что $x_n \to x^*$.
		\item \textbf{Принадлежность шарам:} Для любого фиксированного $N$ все точки $\{x_n\}$ при $n \geqslant N$ лежат в шаре $\overline B_{r_N}(x_N)$. Так как шар является замкнутым множеством, он содержит и свой предел: $x^* \in \overline B_{r_N}(x_N)$. Это верно для любого $N$, следовательно, $x^* \in \bigcap_{n \in \mathbb{N}} \overline B_{r_n}(x_n)$.
		\item \textbf{Единственность:} Если бы существовала вторая точка $y^* \in \bigcap \overline B_{r_n}(x_n)$, то расстояние $\rho(x^*, y^*) \leqslant 2r_n$ для любого $n$. Устремляя $n \to \infty$, получаем $\rho(x^*, y^*) = 0 \Rightarrow x^* = y^*$.
	\end{enumerate}
\end{proof}

\begin{definition}[Нигде не плотное множество]
	Множество $M \subset X$ называется \textit{нигде не плотным} в $X$, если внутренность его замыкания пуста:
	\[ \text{int}(\overline{M}) = \varnothing \]
\end{definition}

\begin{definition}[Критерий «дырявости»]
	Эквивалентно, $M$ \textit{нигде не плотно}, если в любом открытом шаре $B \subset X$ содержится подшар $B_1 \subset B$, не пересекающийся с $M$:
	\[ B_1 \cap M = \varnothing \]
\end{definition}

\begin{definition}[Всюду плотное множество]
	Множество $M \subset X$ называется \textit{всюду плотным} в $X$, если его замыкание совпадает со всем пространством:
	\[ \overline{M} = X \]
	(То есть в любой окрестности любой точки из $X$ есть хотя бы одна точка из $M$).
\end{definition}

\begin{theorem}[Бэра]\label{Ber}
	\textbf{Условия:}
	\begin{itemize}
		\item $X$ --- полное метрическое пространство;
		\item $M_n \subset X$ --- последовательность множеств, нигде не плотных в $X$.
	\end{itemize}
	\textbf{Следовательно:}
	\begin{itemize}
		\item Пространство $X$ нельзя представить в виде счетного объединения таких множеств: $X \neq \bigcup\limits_{n \in \mathbb{N}} M_n$.
	\end{itemize}
\end{theorem}

\begin{proof}
	Предположим от противного: $X = \bigcup_{n=1}^{\infty} M_n$.
	\begin{enumerate}
		\item \textbf{Построение:} 
		Возьмем произвольный шар $\overline{B}_0$. 
		Так как $M_1$ нигде не плотно, внутри открытого шара $B_0$ существует замкнутый подшар $\overline{B}_1$, не пересекающийся с $M_1$. Радиус выберем $r_1 \leqslant 1$.
		
		\item \textbf{Индукция:} 
		Для каждого $n$ внутри открытого шара $B_{n-1}$ найдем замкнутый подшар $\overline{B}_n$ такой, что $\overline{B}_n \cap M_n = \varnothing$ и $r_n \leqslant \frac{1}{2^n}$. (сначала определение, а радиус можем любой взять) Возможность такого выбора гарантирована определением нигде не плотного множества.
		
		\item \textbf{Предел:} 
		Мы получили систему вложенных замкнутых шаров $\overline{B}_n$, радиусы которых стремятся к нулю. В силу \textit{полноты} пространства $X$, существует точка $x^* \in \bigcap_{n=1}^{\infty} \overline{B}_n$.
		
		\item \textbf{Противоречие:} 
		Поскольку $x^* \in \overline{B}_n$ для любого $n$, то $x^* \notin M_n$ ни для одного $n$ (так как $\overline{B}_n \cap M_n = \varnothing$). Следовательно, $x^* \notin \bigcup M_n$. 
		Но по предположению $X = \bigcup M_n$, а значит $x^*$ обязана была попасть в объединение. Противоречие.
	\end{enumerate}
\end{proof}

\begin{idea}{Ловушка для понимания (Почему $\overline{M}_n$?)}
	\textbf{Важный нюанс:} В доказательстве мы строим шары так, чтобы они не пересекались именно с \textit{замыканием} $\overline{M}_n$. 
	
	\textbf{Почему это критично?} Если «уворачиваться» только от самого $M_n$, то предел последовательности шаров $x^*$ может оказаться предельной точкой для $M_n$ (лежать на его границе). Если граница принадлежит $M_n$, то $x^* \in M_n$, и противоречие исчезнет. 
	
	Выбирая шар $\overline{B}_n \cap \overline{M}_n = \varnothing$, мы создаем «зону безопасности», гарантирующую, что $x^*$ не только не совпадет с точками «пыли», но и не окажется на её «краю».
\end{idea}
